% Compile with: pdflatex AA-01_proposed_proof.tex (run twice).
% Standalone source: no external bibliography, figures, or custom fonts needed.
\documentclass[11pt,a4paper]{article}
\usepackage[T1]{fontenc}
\usepackage[utf8]{inputenc}
\usepackage{lmodern}
\usepackage[margin=27mm,headheight=14pt]{geometry}
\usepackage{amsmath,amssymb,amsthm,mathtools}
\usepackage{microtype}
\usepackage{enumitem}
\usepackage{needspace}
\usepackage{fancyhdr}
\usepackage{xurl}
\usepackage[hidelinks]{hyperref}
\hypersetup{
  pdftitle={AA-01: A proposed decision procedure for accurate polynomial evaluation},
  pdfsubject={Self-contained proof candidate for the constant-free finite-tree model},
  pdfauthor={Proof manuscript prepared with ChatGPT}
}
\pagestyle{fancy}
\fancyhf{}
\fancyhead[L]{\small AA-01: proposed proof}
\fancyhead[R]{\small Constant-free finite-tree model}
\fancyfoot[C]{\thepage}
\renewcommand{\headrulewidth}{0.3pt}
\setlength{\parskip}{0.3em}
\setlength{\emergencystretch}{2em}
\setlist[enumerate]{itemsep=0.35em,topsep=0.4em}
\numberwithin{equation}{section}
\newtheorem{theorem}{Theorem}[section]
\newtheorem{lemma}[theorem]{Lemma}
\newtheorem{proposition}[theorem]{Proposition}
\newtheorem{corollary}[theorem]{Corollary}
\theoremstyle{definition}
\newtheorem{definition}[theorem]{Definition}
\newtheorem{example}[theorem]{Example}
\theoremstyle{remark}
\newtheorem{remark}[theorem]{Remark}
\newcommand{\R}{\mathbb R}
\newcommand{\Z}{\mathbb Z}
\newcommand{\N}{\mathbb N}
\newcommand{\NN}{\mathbb N_0}
\newcommand{\Sn}{\mathfrak S_n}
\newcommand{\abs}[1]{\left|#1\right|}
\newcommand{\coefnorm}[1]{\left\lVert#1\right\rVert_{\mathrm{coef}}}
\newcommand{\supp}{\operatorname{supp}}
\newcommand{\rnd}{\operatorname{rnd}}
\newcommand{\sgn}{\operatorname{sgn}}
\newcommand{\YES}{\mathsf{YES}}
\newcommand{\NO}{\mathsf{NO}}
\newcommand{\QE}{\mathsf{QE}}

\title{\textbf{AA-01: A proposed decision procedure\\
for accurate polynomial evaluation}\\[0.5em]
\large Exact real inputs, independently rounded arithmetic,\\
exact comparisons, and no constant inputs}
\author{}
\date{11 September 2026}

\begin{document}
\maketitle
\thispagestyle{plain}

\begin{abstract}
We give a proof candidate for an affirmative answer to AA-01 in its
constant-free, finite-computation-tree formulation.  The proposed
characterization uses the finitely many regions determined by input signs
and the ordering of input absolute values.  After changing to nonnegative
gap coordinates on each region, accurate evaluability is equivalent to a
uniform bound of the coefficientwise absolute polynomial by the absolute
value of the polynomial itself.  Each bound is a first-order sentence over
the reals, so effective real quantifier elimination gives a halting
decision procedure.  The necessity argument couples two nearby executions
while preserving error-dependent branches and all stored rounded values.
The sufficiency argument constructs an evaluator using only exact
sign/order branching and rounded evaluation in the gap coordinates.
\end{abstract}

\noindent\textbf{Verification status.}
This is a research proof candidate developed in the accompanying
conversation.  It has not been independently peer-reviewed or formally
verified.  All model-specific lemmas are proved below; the only external
algorithmic theorem used is effective quantifier elimination over the
reals.  The references supply the problem, background, and that external
theorem, not a prior verification of the proposed characterization.

\section{The precise problem and computation model}\label{sec:model}

Fix an integer $n\ge1$.  The input to the decision problem is the finite
integer coefficient list of a polynomial
\[
 p\in\Z[x_1,\ldots,x_n],\qquad p(0)=0.
\]
The model is the one specified in AA-01~\cite{AA01}.  Background on accurate
evaluation in traditional rounded arithmetic is given in
\cite{DDH06,DDHK08}.

An evaluator is a finite arithmetic computation tree with real-valued
registers and branching control flow.  Its initial numerical quantities
are exactly the input coordinates $x_1,\ldots,x_n$.  Each binary arithmetic
node performs one of $+$, $-$, or multiplication and returns
\begin{equation}\label{eq:rounding}
 (a\circ b)(1+\delta_j),\qquad \circ\in\{+,-,\cdot\}.
\end{equation}
Every execution of a distinct arithmetic node has its own independent
error parameter.  Here ``independent'' imposes no probability distribution:
it means that the errors are separately and arbitrarily selectable within
the prescribed interval.  Equal operands at different nodes need not
produce equal errors.

Negation, copying, storage, and comparisons of available numerical
quantities are exact.  A comparison chooses control flow; its Boolean
outcome is not an additional numerical constant available to arithmetic.
The output is an available numerical quantity.  There are no constant
input nodes and no additional exact arithmetic operations.  In
particular, exact multiplication by an integer is not a primitive.
Previously computed values may be reused without being recomputed.
The tree is finite, including all its branches.  Arithmetic is over the
reals, without overflow or underflow.

Number all rounded nodes in the tree by $1,\ldots,N(P)$.  Errors assigned
to nodes on unused branches are ignored.  Write $P(x,\delta)$ for the
resulting output.

\begin{definition}[Accuracy]\label{def:accuracy}
A tree $P$ accurately evaluates $p$ if $P(x,0)=p(x)$ for all $x\in\R^n$
and
\begin{equation}\label{eq:accuracy}
\begin{gathered}
 \forall\eta\in(0,1)\ \exists u\in(0,1)\quad
 \forall x\in\R^n\ \forall\delta\in[-u,u]^{N(P)}:\\
 \abs{P(x,\delta)-p(x)}\le\eta\abs{p(x)}.
\end{gathered}
\end{equation}
The same tree must work for every $\eta$.  The precision bound $u$ may
depend on $p$, $P$, and $\eta$, but not on $x$.
\end{definition}

The question is whether an ordinary Turing algorithm always halts and
decides whether such a tree exists.  The Turing decision procedure and the
rounded evaluator are different machines: the former may perform exact
symbolic integer computations and invoke a real-algebraic decision
algorithm.

\section{The proposed finite characterization}\label{sec:criterion}

For $\pi\in\Sn$ and $s=(s_1,\ldots,s_n)\in\{-1,1\}^n$, define
$T=T_{\pi,s}:\R^n\to\R^n$ by
\begin{equation}\label{eq:T}
 (Ty)_{\pi(i)}=s_i\sum_{j=i}^n y_j,\qquad 1\le i\le n.
\end{equation}
This is an invertible integer linear map.  Its inverse on its closed
signed-order region is
\begin{equation}\label{eq:gaps}
 y_i=\abs{x_{\pi(i)}}-\abs{x_{\pi(i+1)}}\quad(1\le i<n),
 \qquad y_n=\abs{x_{\pi(n)}}.
\end{equation}
For $y\in(0,\infty)^n$, the image satisfies
\[
 \sgn(x_{\pi(i)})=s_i,
 \qquad \abs{x_{\pi(1)}}>\cdots>\abs{x_{\pi(n)}}>0.
\]
The $2^n n!$ closed regions $T_{\pi,s}[0,\infty)^n$ cover $\R^n$.
For zero coordinates choose either sign, and for ties choose any order.

Expand, after collecting like monomials,
\begin{equation}\label{eq:qW}
 q_T(y)=p(Ty)=\sum_{\alpha\in\NN^n}a_{T,\alpha}y^\alpha,
 \qquad
 W_T(y)=\sum_{\alpha\in\NN^n}\abs{a_{T,\alpha}}y^\alpha.
\end{equation}
The sums are finite, their coefficients are integers, and their constant
terms are zero.  We use the polynomial convention $y_i^0=1$, including
when $y_i=0$.

\Needspace{13\baselineskip}
\begin{theorem}[Signed-order gap criterion]\label{thm:main}
For the model of Section~\ref{sec:model}, the following are equivalent.
\begin{enumerate}[label=\textup{(\roman*)}]
\item\label{item:accurate} The polynomial $p$ has an accurate finite tree.
\item\label{item:dominated} For every $T=T_{\pi,s}$, there is $C_T>0$ such that
\begin{equation}\label{eq:criterion}
 W_T(y)\le C_T\abs{q_T(y)}\qquad\text{for all }y\in[0,\infty)^n.
\end{equation}
\end{enumerate}
When these conditions hold, an accurate evaluator can be constructed whose
only comparisons determine the exact signs and the ordering of the
absolute values of the original inputs.
\end{theorem}

Necessity is proved in Sections~\ref{sec:geometry}--\ref{sec:necessity};
sufficiency is proved in Section~\ref{sec:sufficiency}.  The halting
decision procedure follows in Section~\ref{sec:decision}.
The zero polynomial will be dealt with explicitly and causes no exception
to the criterion.

\section{Geometry of a signed-order region}\label{sec:geometry}

Write $S_i(y)=\sum_{j=i}^n y_j$.  Up to sign and permutation, the raw
inputs are the nested tail sums $S_1(y),\ldots,S_n(y)$.

\begin{lemma}[Relative perturbations of raw quantities]\label{lem:geometry}
Fix $T=T_{\pi,s}$, $0<\varepsilon<1$, and $y,y'\in(0,\infty)^n$ with
\begin{equation}\label{eq:close}
 \abs{y_i'-y_i}\le\varepsilon y_i\qquad(1\le i\le n).
\end{equation}
Let $x=Ty$ and $x'=Ty'$.  Then:
\begin{enumerate}[label=\textup{(\alph*)}]
\item Every signed raw input $r=\sigma x_i$, $\sigma\in\{-1,1\}$,
and its counterpart $r'=\sigma x_i'$ satisfy
\begin{equation}\label{eq:rawperturb}
 \abs{r'-r}\le\varepsilon\abs r.
\end{equation}
\item A sum or difference $L$ of two signed raw inputs is either
identically zero on the region, or is nonzero throughout the open region
and satisfies
\begin{equation}\label{eq:linearperturb}
 \abs{L(x')-L(x)}\le\varepsilon\abs{L(x)}.
\end{equation}
In particular, comparisons between signed raw inputs have the same
outcomes at $x$ and $x'$.
\item A product $a$ of two signed raw inputs and its counterpart $a'$
satisfy
\begin{equation}\label{eq:productperturb}
 \abs{a'-a}\le(2\varepsilon+\varepsilon^2)\abs a
 \le3\varepsilon\abs a.
\end{equation}
\end{enumerate}
\end{lemma}

\begin{proof}
For any $b_i\ge0$, condition~\eqref{eq:close} implies
\[
 \abs{\sum_i b_i y_i'-\sum_i b_i y_i}
 \le\sum_i b_i\abs{y_i'-y_i}
 \le\varepsilon\sum_i b_i y_i.
\]
This proves (a).  A sum or difference of two signed tail sums is either
zero or, up to an overall sign, a nonzero linear form with nonnegative
coefficients: $S_i+S_j$ has this property, and, for $i<j$,
$S_i-S_j=\sum_{k=i}^{j-1}y_k$ does too.  The displayed inequality proves
(b), including preservation of signs and exact equality cases.

For (c), write the two perturbed factors as $r(1+\theta)$ and
$t(1+\phi)$, with $\abs\theta,\abs\phi\le\varepsilon$, using (a).
Their relative product change is
$\theta+\phi+\theta\phi$, whose absolute value is at most
$2\varepsilon+\varepsilon^2$.
\end{proof}

\section{Coupling executions, including error-dependent branches}
\label{sec:coupling}

The next lemma permits each rounded value to be selected away from the
finitely many signed raw inputs.

\begin{lemma}[Avoiding finitely many relative neighborhoods]
\label{lem:avoid}
Let $0<u<1/4$, $n\ge1$, and $\tau=u/(32n)$.  Let $\mathcal R$ be a set
of at most $2n$ nonzero real numbers.  For every $a\in\R$, there is
$\delta\in[-u/4,u/4]$ such that
\begin{equation}\label{eq:separation}
 \abs{a(1+\delta)-r}>\tau\abs r
 \qquad(r\in\mathcal R).
\end{equation}
\end{lemma}

\begin{proof}
If $a=0$, every $\delta$ works because $\tau<1$ and $r\ne0$.
Suppose $a\ne0$, and put $I=[-u/4,u/4]$, an interval of length $u/2$.
For a fixed $r$, its forbidden set is
\[
 F_r=\{\delta\in I:\abs{a(1+\delta)-r}\le\tau\abs r\}.
\]
It is an interval intersected with $I$ and has length at most
$2\tau\abs r/\abs a$.  If it is nonempty, some $\delta\in I$ gives
\[
 (1-\tau)\abs r\le\abs{a(1+\delta)}
 \le(1+u/4)\abs a.
\]
Hence
\[
 \operatorname{length}(F_r)
 \le2\tau\frac{1+u/4}{1-\tau}<3\tau.
\]
For the last inequality, $u/4<1/16$ and $\tau<1/128$, so the factor
multiplying $\tau$ is less than $272/127<3$.
The union of all forbidden intervals has length at most
\[
 6n\tau=\frac{3u}{16}<\frac u2.
\]
It cannot cover $I$.  A point of its complement satisfies every strict
inequality in~\eqref{eq:separation}.
\end{proof}

\begin{lemma}[Deterministic coupling]\label{lem:coupling}
Fix a finite tree $P$ in the prescribed model and $0<u<1/4$.  Set
\begin{equation}\label{eq:parameters}
 \tau=\frac{u}{32n},\qquad
 \varepsilon=\frac{u\tau}{16}=\frac{u^2}{512n}.
\end{equation}
Suppose $x=Ty$, $x'=Ty'$ for one common map $T$, where $y,y'>0$ and
\eqref{eq:close} holds.  There are error assignments
\[
 \delta,\delta'\in[-u,u]^{N(P)}
\]
with the following properties:
\begin{enumerate}[label=\textup{(\roman*)}]
\item the two executions follow exactly the same control-flow path;
\item every rounded node on that path produces the same numerical value
in the two executions;
\item if their outputs are $z=P(x,\delta)$ and $z'=P(x',\delta')$, then
\begin{equation}\label{eq:outputcoupling}
 \abs{z'}\le(1+\varepsilon)\abs z.
\end{equation}
\end{enumerate}
The parameter $\varepsilon$ depends only on $n$ and $u$, not on the size
of $P$.
\end{lemma}

\begin{proof}
Every available numerical quantity has one of two provenance types:
it is an exact signed copy of an original input coordinate, or an exact
signed copy of a value previously produced by rounded arithmetic.
Copies and negations preserve these types.  They are provenance labels,
not claims that differently labelled values can never coincide.

We construct the two executions simultaneously, assigning an error when
a rounded node is reached.  For every rounded value $v$ produced in the
execution on $x$, maintain
\begin{equation}\label{eq:invariant}
 \abs{v-r}>\tau\abs r
 \qquad\text{for all }r\in\mathcal R_x
 :=\{x_1,-x_1,\ldots,x_n,-x_n\}.
\end{equation}
All targets are nonzero because $y>0$.  The same separation holds for
$-v$, since $\mathcal R_x$ is symmetric under negation.  We also maintain
that the counterpart of every such rounded value is numerically identical
in the second execution.  Signed raw inputs instead have the naturally
perturbed counterparts from Lemma~\ref{lem:geometry}.

Consider corresponding arithmetic nodes, with exact results $a,a'$
before the new rounding.  The following cases exhaust all their operands.

\smallskip
\noindent\emph{Both operands have rounded provenance.}
The operands match exactly, so $a'=a$, including when $a=0$.

\smallskip
\noindent\emph{Both operands have raw-input provenance.}
For addition or subtraction, Lemma~\ref{lem:geometry}(b) gives either
$a=a'=0$ or a relative change bounded by $\varepsilon$.
For multiplication, both factors are nonzero and part (c) bounds the
relative change by $3\varepsilon$.

\smallskip
\noindent\emph{One operand has each provenance type, with multiplication.}
Write them as a matched value $v$ and a signed raw input $r$.
If $v=0$, then $a=a'=0$.  Otherwise,~\eqref{eq:rawperturb} bounds the
relative change by $\varepsilon$.

\smallskip
\noindent\emph{One operand has each provenance type, with addition or
subtraction.}
Up to an overall sign, $a=v\pm r$, where $v$ is a signed copy of a prior
rounded value.  The invariant gives
\[
 \abs a>\tau\abs r,
 \qquad
 \abs{a'-a}\le\varepsilon\abs r.
\]
In particular, $a\ne0$, and
\begin{equation}\label{eq:mixedchange}
 \frac{\abs{a'-a}}{\abs a}
 <\frac{\varepsilon}{\tau}=\frac u{16}.
\end{equation}

Thus every zero case has $a'=0$, and every nonzero case has
\begin{equation}\label{eq:centerchange}
 a'=a(1+\theta),\qquad\abs\theta\le u/16.
\end{equation}
Here the raw-input cases use $3\varepsilon\le u/16$, which follows from
\eqref{eq:parameters} and $u<1/4$, $n\ge1$.

By Lemma~\ref{lem:avoid}, choose the new error $\delta_j\in[-u/4,u/4]$
so that the first execution's new value
$v_{\mathrm{new}}=a(1+\delta_j)$ satisfies~\eqref{eq:invariant}.
If $a=0$, choose $\delta_j'=0$ in the second execution; both new values
are zero.  If $a\ne0$, set
\begin{equation}\label{eq:matcherror}
 \delta_j'=\frac{1+\delta_j}{1+\theta}-1.
\end{equation}
Then
\[
 a'(1+\delta_j')=a(1+\delta_j),
 \qquad
 \abs{\delta_j'}
 \le\frac{u/4+u/16}{1-u/16}
 =\frac{5u}{16-u}<u.
\]
The new rounded values match and the invariant is preserved.

It remains to check each comparison, since its operands may depend on
rounding errors.  Two quantities of rounded provenance match in both
executions, so their comparison is identical, including equality.
Comparisons between signed raw inputs are preserved by
Lemma~\ref{lem:geometry}.  For a mixed comparison, its difference $v-r$
has absolute value greater than $\tau\abs r$, whereas the perturbation
of $r$ has absolute value at most $\varepsilon\abs r$.  Since
$\varepsilon<\tau$, its strict sign cannot change, and equality cannot
occur in either execution.  This also covers signed copies of rounded
values by symmetry of the target set.

Starting with the input registers, these arguments inductively preserve
the path through exact comparisons, negations, copies, and rounded nodes.
They do not recompute stored quantities or assign new errors to a reuse.
The path is finite.  The chosen errors can therefore be extended to full
vectors for the two finite trees by assigning zero to every unused node.
The independent-error convention makes these legal assignments, even
though they were constructed successively.

The common output location either has rounded provenance, in which case
$z'=z$, or is a signed raw-input selection, in which case
\eqref{eq:rawperturb} gives~\eqref{eq:outputcoupling}.
The argument also includes paths with no rounded nodes.
\end{proof}

\begin{remark}\label{rem:no-leaf}
The lemma does not assert that an individual leaf's zero-error expression
equals $p$ on all inputs.  It compares two actual executions with their
selected errors.  Nor does it assume that a leaf is reachable at zero
error.  These distinctions are essential when branching depends on
rounded intermediate values.
\end{remark}

\section{Necessity of the criterion}\label{sec:necessity}

\begin{proposition}[Local control under gap perturbations]
\label{prop:local}
If $P$ accurately evaluates $p$, then there is $\varepsilon>0$ such that,
for every $T$ and all $y,y'>0$ satisfying~\eqref{eq:close},
\begin{equation}\label{eq:local}
 \abs{p(Ty')}\le2\abs{p(Ty)}.
\end{equation}
The same $\varepsilon$ works for every $T$ and every $y$.
\end{proposition}

\begin{proof}
Use~\eqref{eq:accuracy} with $\eta=1/4$.  By shrinking its precision
bound, take $0<u<1/4$.  Use the $\varepsilon$ in
Lemma~\ref{lem:coupling}.  For the coupled outputs $z,z'$,
\[
 \abs{z-p(Ty)}\le\tfrac14\abs{p(Ty)},
 \qquad
 \abs{z'-p(Ty')}\le\tfrac14\abs{p(Ty')}.
\]
Therefore
\begin{align*}
 \abs{p(Ty')}
 &\le\tfrac43\abs{z'}
 \le\tfrac43(1+\varepsilon)\abs z\\
 &\le\tfrac53(1+\varepsilon)\abs{p(Ty)}
 \le2\abs{p(Ty)}.
\end{align*}
The last inequality follows, for example, from
$\varepsilon=u^2/(512n)<1/8192<1/5$.
No division by $p(Ty)$ is used, so zeros are included.
\end{proof}

We next give an explicit interpolation proof of the finite-dimensional
coefficient bound needed below.

\begin{lemma}[Coefficient extraction on a fixed box]\label{lem:coefficients}
Fix $n\ge1$, an integer $d\ge0$, and $0<\varepsilon<1$.  There is a
finite constant $K=K(n,d,\varepsilon)>0$ such that every real polynomial
$R(t_1,\ldots,t_n)$ of total degree at most $d$ satisfies
\begin{equation}\label{eq:coefbound}
 \coefnorm R
 :=\sum_\alpha\abs{\operatorname{coeff}_\alpha R}
 \le K\sup_{t\in[1-\varepsilon,1+\varepsilon]^n}\abs{R(t)}.
\end{equation}
\end{lemma}

\begin{proof}
For $d=0$, take $K=1$.  For $d\ge1$, choose $d+1$ distinct points
$t_0,\ldots,t_d$ in $[1-\varepsilon,1+\varepsilon]$ and define the
univariate Lagrange polynomials
\[
 L_j(t)=\prod_{\substack{0\le k\le d\\k\ne j}}
 \frac{t-t_k}{t_j-t_k}.
\]
Successive univariate interpolation in each coordinate gives
\[
 R(t)=\sum_{j_1=0}^d\cdots\sum_{j_n=0}^d
 R(t_{j_1},\ldots,t_{j_n})\prod_{i=1}^n L_{j_i}(t_i).
\]
This applies because the degree in each separate variable is at most
$d$.  The coefficient norm is subadditive.  Because the factors in the
product involve distinct variables, the coefficient norm of that
product is $\prod_i\coefnorm{L_{j_i}}$.  Thus~\eqref{eq:coefbound} holds
with the finite positive constant
\[
 K=\left(\sum_{j=0}^d\coefnorm{L_j}\right)^n.
\]
\end{proof}

\begin{proof}[Proof of Theorem~\ref{thm:main}, necessity]
Suppose $p$ has an accurate tree.  If $p=0$, then $q_T=W_T=0$ for every
$T$, and the criterion is immediate.  Otherwise let $d=\deg p$ and use
$\varepsilon$ from Proposition~\ref{prop:local}.

Fix $T$ and $y\in(0,\infty)^n$.  For
$t\in[1-\varepsilon,1+\varepsilon]^n$, the point
$y'=t\odot y=(t_1y_1,\ldots,t_ny_n)$ satisfies~\eqref{eq:close}.
Consequently,
\begin{equation}\label{eq:boxbound}
 \sup_{t\in[1-\varepsilon,1+\varepsilon]^n}
 \abs{q_T(t\odot y)}\le2\abs{q_T(y)}.
\end{equation}
Apply Lemma~\ref{lem:coefficients} to the polynomial in $t$
\[
 R_y(t)=q_T(t\odot y)
 =\sum_\alpha a_{T,\alpha}y^\alpha t^\alpha.
\]
Since $y>0$, its coefficient norm is exactly $W_T(y)$.
Equation~\eqref{eq:boxbound} yields
\[
 W_T(y)\le2K\abs{q_T(y)}.
\]
Both sides are continuous functions of $y$, so the same inequality holds
on the closure $[0,\infty)^n$.  This proves~\eqref{eq:criterion}, with
$C_T=2K$, for each $T$.
\end{proof}

\section{Constructing an accurate evaluator}\label{sec:sufficiency}

For completeness, the following evaluation bound explicitly retains all
error factors introduced by stored-value reuse.

\begin{lemma}[Termwise evaluation with relative input errors]
\label{lem:termwise}
Let
\[
 q(y)=\sum_{\alpha\ne0}a_\alpha y^\alpha\in\Z[y_1,\ldots,y_n]
\]
be nonzero with zero constant term.  Suppose the available nonnegative
data have the form $\widehat y_i=y_i(1+\xi_i)$, with $y_i\ge0$ and
$\abs{\xi_i}\le u<1$.  There is a fixed constant-free arithmetic
program for $q$, using these data, with output $\widehat q$ and an integer
$L\ge0$ such that
\begin{equation}\label{eq:termwise}
 \abs{\widehat q-q(y)}
 \le\bigl((1+u)^L-1\bigr)
 \sum_{\alpha\ne0}\abs{a_\alpha}y^\alpha.
\end{equation}
Here the same bound $u$ applies to all its new arithmetic errors, and
$L$ depends only on the chosen program and $q$.
\end{lemma}

\begin{proof}
For every nonzero coefficient, form $y^\alpha$ by multiplying its
$\abs\alpha=\sum_i\alpha_i\ge1$ factors, beginning with its first
factor rather than with a constant $1$.  To implement its integer
coefficient, take $\abs{a_\alpha}$ occurrences of this computed monomial
in a sum and use exact negation if $a_\alpha<0$.  Then add all these
signed occurrences, starting with the first one rather than a constant
$0$.  This is a finite program; large coefficients affect its size, not
its existence.  Reuse of the computed monomial is permitted.

Substitute each input factor $y_i(1+\xi_i)$ into this program and expand
only the additions, including their rounding factors.  If $\zeta$ lists
all input-error and arithmetic-error variables, its output has the form
\begin{equation}\label{eq:occurrences}
 \widehat q
 =\sum_{r=1}^{M}\sigma_r y^{\alpha(r)}
       \prod_j(1+\zeta_j)^{e_{rj}},
 \qquad \sigma_r\in\{-1,1\},\quad e_{rj}\in\NN.
\end{equation}
There are exactly $\abs{a_\alpha}$ occurrences with exponent $\alpha$,
and all these occurrences have sign $\sgn(a_\alpha)$.  In particular,
\[
 \sum_{r=1}^M\sigma_r y^{\alpha(r)}=q(y),
 \qquad
 \sum_{r=1}^M y^{\alpha(r)}
 =\sum_\alpha\abs{a_\alpha}y^\alpha.
\]
The exponents $e_{rj}$ retain multiplicities of shared input errors and
shared stored values; no independent error is assigned to a reuse.
The expansion~\eqref{eq:occurrences} follows directly because no
multiplication of two sums is required in the chosen program.

Let $L=\max_r\sum_j e_{rj}$.  When $\abs{\zeta_j}\le u<1$, each error
product lies between $(1-u)^L$ and $(1+u)^L$.  Furthermore,
\[
 1-(1-u)^L\le(1+u)^L-1,
\]
by the binomial theorem (also when $L=0$).  Thus every error product in
\eqref{eq:occurrences} differs from $1$ by at most $(1+u)^L-1$.
Subtract $q(y)$ and use the triangle inequality to obtain
\eqref{eq:termwise}.
\end{proof}

\begin{proof}[Proof of Theorem~\ref{thm:main}, sufficiency]
Assume~\eqref{eq:criterion}.  If $p=0$, the single rounded subtraction
$x_1-x_1$ returns zero for every error and is an accurate evaluator.
Suppose $p\ne0$.

First determine the sign of each $x_i$ by exactly comparing $x_i$ with
$-x_i$.  Assign sign $+1$ when $x_i=0$.  Form the absolute values by
copying or exact negation, and sort them by exact comparisons; break ties
by a fixed index rule.  For fixed $n$, this can be implemented by a finite
comparison tree.  No arithmetic constant is used: the signs in this
description specify copying or negation, not multiplication by a stored
number $1$ or $-1$.

This selects $T=T_{\pi,s}$ with $x=Ty$ and $y\ge0$ as
in~\eqref{eq:gaps}.  Compute the first $n-1$ gaps by rounded subtraction
of the exact signed raw inputs.  Keep the smallest absolute input as the
last gap.  Hence the available values are
\[
 \widehat y_i=y_i(1+\xi_i)\quad(i<n),
 \qquad \widehat y_n=y_n,
\]
where each $\xi_i$ is its own rounded-node error.  In particular, zero
gaps are exactly zero.

On this branch, apply Lemma~\ref{lem:termwise} to $q_T$, treating the
last input error as zero.  For a finite branch-dependent integer $L_T$,
\[
 \abs{P(x,\delta)-p(x)}
 =\abs{\widehat q_T-q_T(y)}
 \le\bigl((1+u)^{L_T}-1\bigr)W_T(y)
 \le C_T\bigl((1+u)^{L_T}-1\bigr)\abs{p(x)}.
\]
There are only finitely many branches.  Set
\[
 C_* =\max_T C_T,\qquad L_* =\max_T L_T.
\]
For every $\eta\in(0,1)$, there is $u\in(0,1)$ such that
$C_*((1+u)^{L_*}-1)\le\eta$, by continuity at $u=0$.
If $L_*=0$, the bound is already zero.
This proves~\eqref{eq:accuracy} uniformly in $x$ and all allowed errors.

At $\delta=0$, the gap computations and subsequent polynomial
computations are exact, so $P(x,0)=p(x)$.  At any zero of $p$, the same
bound forces the output to be zero for every allowed error.  Equivalently,
$W_T(y)=0$ there, so every contributing monomial has a zero gap factor
and remains zero under the prescribed arithmetic.

The construction, including the branch tests and the finite arithmetic
on each branch, depends only on $p$.  It does not depend on $\eta$, $u$,
or the values of the constants $C_T$.  Thus it is one fixed tree valid
for every requested accuracy, as required.
\end{proof}

\section{The halting decision procedure}\label{sec:decision}

We invoke the standard effective form of real quantifier elimination:
there is a Turing algorithm that decides every first-order sentence over
$\R$ whose atomic formulas are polynomial equalities and inequalities
with rational coefficients.  For a primary reference giving a complete
formally verified quantifier-elimination algorithm, see~\cite{KTP23}.
That verification concerns the external quantifier-elimination algorithm,
not the new coupling argument in this manuscript.

\begin{corollary}[Proposed affirmative answer to AA-01]
\label{cor:decision}
There is an always-halting Turing algorithm which, from the integer
coefficient list of $p$ with $p(0)=0$, decides whether an accurate finite
tree exists in the stated model.  On a positive instance it can also
output the finite evaluator constructed in Section~\ref{sec:sufficiency}.
\end{corollary}

\begin{proof}
Use the following algorithm.
\begin{enumerate}[label=\arabic*.]
\item If $p=0$, return $\YES$ and, when requested, the tree that computes
$x_1-x_1$.
\item Enumerate the finite set of pairs
$(\pi,s)\in\Sn\times\{-1,1\}^n$.
For each pair, symbolically substitute~\eqref{eq:T} into $p$, expand and
collect its integer coefficients, and construct $q_T,W_T$
from~\eqref{eq:qW}.
\item For each pair decide the closed first-order sentence
\begin{equation}\label{eq:QE}
 \exists C\;\Bigl(C>0\ \wedge\
 \forall y_1\cdots\forall y_n\;
 \Bigl[
  \bigl(\bigwedge_{i=1}^n y_i\ge0\bigr)
  \Longrightarrow
  W_T(y)^2\le C^2 q_T(y)^2
 \Bigr]\Bigr).
\end{equation}
If any sentence is false, return $\NO$.
\item Otherwise return $\YES$ and construct the sign/order tree and its
branch evaluators as in Section~\ref{sec:sufficiency}.
\end{enumerate}
All symbolic operations are finite operations on integer coefficient
lists.  Sentence~\eqref{eq:QE} involves only integer-coefficient
polynomials in $C,y_1,\ldots,y_n$, so its decision terminates.  There are
only $2^n n!$ such sentences.  The whole algorithm therefore terminates.

For $y\ge0$ we have $W_T(y)\ge0$.  Since $C>0$, the squared inequality
in~\eqref{eq:QE} is equivalent to~\eqref{eq:criterion}; this includes
points where $q_T(y)=0$.  Correctness follows in both directions from
Theorem~\ref{thm:main}.

Neither the decision procedure nor the evaluator construction requires
knowing the size of a putative accurate tree.  Also, the quantified
constant $C$ is used only in the symbolic decision test and the analysis;
it is not a numerical constant input to the evaluator.
\end{proof}

\section{Boundary cases, examples, and scope}\label{sec:scope}

\Needspace{6\baselineskip}
\begin{example}[An allowable input difference]
For $p(x)=x_i-x_j$, each $q_T$ is either zero or, up to an overall sign,
a linear form with nonnegative coefficients.  Thus
$W_T(y)=\abs{q_T(y)}$ on $y\ge0$ and the criterion accepts $p$.
The single rounded subtraction is already an accurate evaluator.
\end{example}

\begin{example}[Three-input summation]
Take $p=x_1+x_2+x_3$ and the signed-order substitution
\[
 x_1=y_1+y_2+y_3,\qquad
 x_2=-(y_2+y_3),\qquad x_3=-y_3.
\]
Then $q_T=y_1-y_3$ and $W_T=y_1+y_3$.  At $y=(1,1,1)$, the former
vanishes while the latter is $2$.  Sentence~\eqref{eq:QE} is therefore
false for this region, so the criterion rejects the polynomial.  This
agrees with the impossibility example discussed in~\cite{DDH06,DDHK08}.
\end{example}

\begin{example}[Bounded cancellation is permitted]
Let $p(x)=x^5-x^3+x$.  In either one-variable sign region,
\[
 W(y)=y^5+y^3+y,\qquad
 \abs{q(y)}=y(y^4-y^2+1),\qquad y\ge0.
\]
The factor $y^4-y^2+1$ is strictly positive, because it equals
$(y^2-1/2)^2+3/4$.  Moreover,
\[
 3\abs{q(y)}-W(y)=2y(y^2-1)^2\ge0.
\]
Thus $C=3$ works, even though the polynomial has a negative coefficient.
The fractions in the positivity identity belong only to this proof, not
to the evaluator's constant inputs.
\end{example}

\paragraph{Closed regions and zeros.}
Necessity is first proved on strictly ordered, nonzero inputs, then
extended by continuity to the entire nonnegative orthant of gap
coordinates.  Sufficiency uses those closed-region inequalities, so it
covers tied magnitudes, zero coordinates, the origin, and zeros of $p$.
No omission of a lower-dimensional set is made.

\paragraph{Stored values and error-dependent comparisons.}
The necessity proof couples actual register values and explicitly checks
every provenance type of comparison.  In particular, it does not turn a
stored-value program into a formula with newly independent errors at
reused occurrences.  The formal expansion used for the deliberately
constructed evaluator in Lemma~\ref{lem:termwise} retains shared error
variables and their multiplicities.

\paragraph{Limits of the stated result.}
The theorem concerns whole-space polynomial evaluation in the exact model
of Section~\ref{sec:model}.  It does not assert a characterization for
arbitrary restricted domains, additional arithmetic primitives, nonzero
constant inputs, numerical access to comparison outcomes, deterministic
rounding rules tying errors at repeated operands, or bit-level
floating-point algorithms.  Those changes alter premises of the coupling
argument.  No complexity or practical-efficiency bound is claimed.

\paragraph{Logical dependencies.}
The substantive new step is Lemma~\ref{lem:coupling}.  It implies
Proposition~\ref{prop:local}; polynomial interpolation then proves
necessity.  A direct termwise construction proves sufficiency.  The
standard real-quantifier-elimination theorem yields the final Turing
decision procedure.  The manuscript supplies all these model-specific
steps but remains a proposed proof pending independent verification.

\Needspace{22\baselineskip}
\begin{thebibliography}{9}

\bibitem{AA01}
A.~Townsend,
\emph{Open Problems in Numerical Linear Algebra},
AA-01: ``Deciding accurate evaluability of real polynomials,''
statement in the MColbrook repository used in the accompanying
conversation.
\url{https://github.com/MColbrook/OpenProblemsInNLA/blob/main/arithmetic-and-complexity/AA-01/README.md}.

\bibitem{DDH06}
J.~Demmel, I.~Dumitriu, and O.~Holtz,
\emph{Toward accurate polynomial evaluation in rounded arithmetic},
in \emph{Foundations of Computational Mathematics: Santander 2005},
Cambridge University Press, 2006, pp.~36--105.
ArXiv version 2, 18 January 2006:
\url{https://arxiv.org/abs/math/0508350v2}.

\bibitem{DDHK08}
J.~Demmel, I.~Dumitriu, O.~Holtz, and P.~Koev,
\emph{Accurate and efficient expression evaluation and linear algebra},
Acta Numerica \textbf{17} (2008), 87--145.
Definitions~3.3--3.4 and Section~3.3, especially Section~3.3.7.
\url{https://doi.org/10.1017/S0962492906350015}.

\bibitem{KTP23}
K.~Kosaian, Y.~K.~Tan, and A.~Platzer,
\emph{A First Complete Algorithm for Real Quantifier Elimination in
Isabelle/HOL},
CPP 2023.
\url{https://doi.org/10.1145/3573105.3575672}.
Preprint:
\url{https://arxiv.org/abs/2209.10978v2}.

\end{thebibliography}
\end{document}
